\documentclass[12pt]{article}
\usepackage{amsmath, amssymb, amsthm}
\usepackage{geometry}
\geometry{margin=1in}

\title{Tutorial--1: Solution\\
CSE 400: Fundamentals of Probability in Computing\\
School of Engineering and Applied Science (SEAS), Ahmedabad University}
\date{February 13, 2026}

\begin{document}

\maketitle

\section*{Definitions and Notations}

\subsection*{Poisson Probability Mass Function}
\[
P(X = k) = \frac{e^{-\lambda} \lambda^k}{k!}
\]

\subsection*{Gaussian PDF}
\[
f_X(x) = \frac{1}{\sqrt{2\pi\sigma^2}} 
\exp\left(-\frac{(x-m)^2}{2\sigma^2}\right)
\]

\subsection*{Gaussian Relations}
\[
F_X(x) = \Phi\left(\frac{x-m}{\sigma}\right),
\quad
\Pr(X > x) = Q\left(\frac{x-m}{\sigma}\right)
\]

\subsection*{Cumulative Distribution Function}
\[
F(a) = \sum_{x \le a} p(x)
\]

If $X$ is discrete with values $x_1 < x_2 < x_3 < \cdots$, then $F$ is a step function.

\section*{Assumptions and Preliminaries}

\begin{itemize}
\item Multinomial counting formula
\item Uniform distribution over an interval
\item Complement rule
\item De Morgan's law
\item Conditional probability formula
\item Binomial distribution
\item Poisson distribution
\item Gaussian distribution
\item Independence of events
\item Jury decision conditioning on guilt with probability $\alpha$
\end{itemize}

\section*{Main Results}

\subsection*{Q1}

\textbf{(a) Total number of ways}
\[
\frac{20!}{(5!)^4 4!}
\]

\textbf{(b) Probability all platters same cuisine}
\[
P = \frac{(5!)^4 4!}{20!}
\]

\subsection*{Q2}

\[
P(\text{wait} < 5) = \frac{10}{30} = \frac{1}{3}
\]

\[
P(\text{wait} > 10) = \frac{10}{30} = \frac{1}{3}
\]

\subsection*{Q3}

\[
P(L)=0.15,\quad P(E)=0.25,\quad P(L \cap E)=0.08
\]

\[
P(L^c \cap E^c)=0.68,\quad P(E^c)=0.75
\]

\[
P(L^c \mid E^c) = \frac{0.68}{0.75} \approx 0.907
\]

\subsection*{Q4}

\[
P(X=0)=e^{-0.2}=0.8187
\]

\[
P(X\ge2)=1-0.8187-0.1637 \approx 0.0175
\]

\subsection*{Q5}

\[
P(X\ge2)=0.8641
\]

\[
P(X\le1)=0.1359
\]

\subsection*{Q6}

\[
P(X\ge3)=0.5768
\]

\[
P(X\ge3 \mid X\ge1)=\frac{0.5768}{0.9502} \approx 0.607
\]

\subsection*{Q7}

\[
P(X \le 0) = 1 - Q\left(\frac{3}{2}\right)
\]

\[
P(X > 4) = Q\left(\frac{7}{2}\right)
\]

\[
P(|X+3|<2) = 1 - 2Q(1)
\]

\[
P(|X-2|>1) = 1 - Q(2) + Q(3)
\]

\subsection*{Q8}

\[
\alpha \sum_{i=8}^{12} 
\binom{12}{i} \theta^i (1-\theta)^{12-i}
+
(1-\alpha)
\sum_{i=5}^{12}
\binom{12}{i} \theta^i (1-\theta)^{12-i}
\]

\subsection*{Q9}

Since
\[
p(i)=c\frac{\lambda^i}{i!}
\]
we obtain
\[
c=e^{-\lambda}
\]

\[
P(X=0)=e^{-\lambda}
\]

\[
P(X>2)=1-e^{-\lambda}-\lambda e^{-\lambda}-\frac{\lambda^2 e^{-\lambda}}{2}
\]

\subsection*{Step Function Example}

\[
F(a)=
\begin{cases}
0, & a<1 \\
\frac14, & 1\le a<2 \\
\frac34, & 2\le a<3 \\
\frac78, & 3\le a<4 \\
1, & a\ge4
\end{cases}
\]

\subsection*{Q10}

\textbf{(a)}
\[
3(p-1)^2(2p-1)>0
\quad \Longleftrightarrow \quad
p>\frac12
\]

\textbf{(b)}
\[
P_{2k+1}-P_{2k-1}
=
\binom{2k-1}{k}
p^k(1-p)^k(2p-1)
\]

\[
>0
\quad \Longleftrightarrow \quad
p>\frac12
\]

\section*{Conclusion (Working Answers)}

\begin{itemize}
\item Q1: $\frac{20!}{(5!)^4 4!}$, \quad $\frac{(5!)^4 4!}{20!}$
\item Q2: $\frac13, \frac13$
\item Q3: $0.907$
\item Q4: $0.8187, 0.0175$
\item Q5: $0.8641, 0.1359$
\item Q6: $0.5768, 0.607$
\item Q7: $1-Q\left(\frac32\right), Q\left(\frac72\right), 1-2Q(1), 1-Q(2)+Q(3)$
\item Q8: Full conditional binomial expression
\item Q9: $e^{-\lambda}, 1-e^{-\lambda}-\lambda e^{-\lambda}-\frac{\lambda^2 e^{-\lambda}}{2}$
\item Q10: $p>\frac12$
\end{itemize}

\end{document}